% Bagian: sets-functions-relations
% Bab: arithmetization
% Subbagian: bilangan rasional

\documentclass[../../../include/open-logic-section]{subfiles}

\begin{document}

\olfileid[id]{sfr}{arith}{rat}
\olsection{Dari $\Int$ ke $\Rat$}

Kita baru saja melihat cara membangun bilangan bulat dari bilangan asli,
dengan menggunakan teori himpunan naif. Sekarang kita akan melihat cara
membangun bilangan rasional dari bilangan bulat dengan cara yang sangat
mirip. Pengamatan awal kita adalah:
\begin{enumerate}
	\item Setiap bilangan rasional dapat ditulis dalam bentuk $\nicefrac{i}{j}$,
	dengan $i$ dan $j$ keduanya bilangan bulat, tetapi $j$ taknol.
	\item Informasi yang dikodekan dalam suatu ekspresi $\nicefrac{i}{j}$
	dapat pula dikodekan dalam pasangan terurut $\tuple{ i, j}$.
\end{enumerate}
Pendekatan yang jelas ialah memandang bilangan rasional \emph{sebagai}
pasangan terurut yang diambil dari $\Int \times (\Int \setminus \{0_\Int\})$.
Namun, seperti sebelumnya, pendekatan itu agak terlalu naif, karena kita
menginginkan $\nicefrac{3}{2} = \nicefrac{6}{4}$, sedangkan $\tuple{ 3,
2}\neq \tuple{ 6, 4}$. Secara lebih umum, kita menginginkan:
\[
	\nicefrac{a}{b} = \nicefrac{c}{d} \text{ jika dan hanya jika } a \times d = b \times c
\]
Untuk memperolehnya, kita mendefinisikan suatu {relasi ekuivalensi} pada
$\Int \times (\Int \setminus \{0_\Int\})$ sebagai berikut:
\[
	\tuple{ a, b }\Ratequiv \tuple{ c, d} \text{ jika dan hanya jika }a \times d = b \times c
\]
Kita harus memeriksa bahwa relasi ini merupakan relasi ekuivalensi. Hal ini
sangat mirip dengan kasus $\Intequiv$, dan kita serahkan sebagai latihan.
\begin{prob}
Tunjukkan bahwa $\Ratequiv$ merupakan relasi ekuivalensi.
\end{prob}
Namun, relasi ini memungkinkan kita mengatakan:
\begin{defn}
Bilangan rasional adalah kelas-kelas ekuivalensi, di bawah $\Ratequiv$, dari
pasangan-pasangan bilangan bulat (dengan elemen kedua yang taknol). Yakni,
$\Rat =
\equivclass{(\Int \times (\Int\setminus \{0_\Int\}))}{\Ratequiv}$.
\end{defn}

Seperti halnya bilangan bulat, kita juga ingin mendefinisikan beberapa
operasi dasar. Jika $\equivrep{i,j}{\Ratequiv}$ adalah kelas ekuivalensi di
bawah $\Ratequiv$ yang memuat $\tuple{i, j}$ sebagai !!a{element}, kita
menetapkan:
\begin{align*}
	\equivrep{a, b}{\Ratequiv} + \equivrep{c, d}{\Ratequiv} &= \equivrep{ad + bc,  bd}{\Ratequiv}\\
	\equivrep{a, b}{\Ratequiv} \times \equivrep{c, d}{\Ratequiv} &= \equivrep{a  c, b d}{\Ratequiv}.
\intertext{Untuk mendefinisikan $r \leq s$ pada bilangan rasional ini, kita menggunakan fakta bahwa
$r \le s$ jika dan hanya jika $s - r$ tidak negatif, yaitu $s - r$ dapat ditulis sebagai
$\nicefrac{i}{j}$ dengan $i$ taknegatif dan $j$~positif:}
	\equivrep{a, b}{\Ratequiv} \leq \equivrep{c, d}{\Ratequiv} &\text{ jika dan hanya jika }
	\equivrep{c, d}{\Ratequiv} - \equivrep{a, b}{\Ratequiv} =
	\equivrep{i_\Int, j_\Int}{\Ratequiv}
\end{align*}
untuk suatu $i \in \Nat$ dan $0 \neq j \in \Nat$.

Selanjutnya, kita perlu memeriksa bahwa definisi-definisi tersebut
berperilaku sebagaimana \emph{semestinya}; hal ini kita tempatkan dalam
\olref[check]{sec}. Namun, semuanya memang berperilaku demikian!{} Terakhir, kita
menginginkan suatu cara untuk memperlakukan bilangan bulat \emph{sebagai}
bilangan rasional; maka, untuk setiap $i \in \Int$, kita menetapkan bahwa
$i_\Rat = \equivrep{i, 1_\Int}{\Ratequiv}$. Sekali lagi, kita memeriksa
dalam \olref[check]{sec} bahwa semua ini berperilaku dengan tepat.

\begin{prob}
Tunjukkan bahwa $(i + j)_\Rat = i_\Rat+ j_\Rat$ dan $(i \times j)_\Rat =
i_\Rat \times j_\Rat$ dan $i \leq j \liff i_\Rat \leq j_\Rat$, untuk
sembarang $i, j \in \Int$.
\end{prob}

\end{document}
